\section{ЗАКЛЮЧЕНИЕ}

В работе был предложен новый метод предсказания микроРНК, ассоциированных с какими-либо процессами или заболеваниями, по сопоставлению топологических характеристик сетей сигнальных путей, участвующих в процессе, с сетью генов-мишеней, рассматриваемой микроРНК.

Рассматривались микроРНК гены-мишени которых экспрессируются в ткани миокарда и сигнальные пути ассоциированные с её репарацией и восстановлением.

В результате, как перспективные для дальнейшего изучения выделены микроРНК:

Для сигнального пути HIF-1a: miR651-5p, miR-147a; для VEGF: miR-374b-5p и miR-369-3p; для FGF2: miR-3652, miR-4430 и miR-6134; для TNF-a: miR-23a-3p; для  IGF-1, IL2, PDGF: miR-2117, miR-196b-3p и miR-155-3p, соответственно.